% =========================================================================
% Animal Care & Research Ethics Compliance Statement
% =========================================================================

\chapter*{Statement of Research Ethics and Animal Care}
\addcontentsline{toc}{chapter}{Statement of Research Ethics and Animal Care}

All research protocols, field capture procedures, animal handling, and laboratory experiments conducted during this dissertation were reviewed and formally approved by the Institutional Animal Care and Use Committee (IACUC) at the host institution under approved protocol:

\begin{center}
\fbox{\bfseries IACUC Protocol Authorization Number: \texttt{IACUC-2023-BIO-0492-A}}
\end{center}

\section*{Ethical Standards and Regulatory Compliance}
Field and laboratory investigations adhered strictly to the \textit{Guidelines for Use of Live Amphibians and Reptiles in Field and Laboratory Research} formulated jointly by the American Society of Ichthyologists and Herpetologists (ASIH), The Herpetologists' League (HL), and the Society for the Study of Amphibians and Reptiles (SSAR), as well as the National Research Council's \textit{Guide for the Care and Use of Laboratory Animals} \parencite{nrc2011animalguide}.

\section*{Permits and International Authorizations}
Specimen collection, non-invasive genetic sampling, and exportation of biological materials were carried out under valid scientific research and collection permits issued by sovereign national wildlife authorities:
\begin{itemize}
  \item \textbf{Republic of Peru:} Servicio Nacional Forestal y de Fauna Silvestre (SERFOR), Scientific Research Authorization Permit No.~\texttt{AUT-IFS-2023-088}; Exportation Permit No.~\texttt{EXP-2023-SERFOR-0412}; CITES Export Permit No.~\texttt{PE-CITES-2023-01948}.
  \item \textbf{Republic of Ecuador:} Ministerio del Ambiente, Agua y Transici{\'o}n Ecol{\'o}gica (MAATE), Framework Scientific Research Authorization No.~\texttt{MAATE-DBI-CM-2023-0115}; Genetic Resource Access Agreement No.~\texttt{MAATE-RGN-2023-044}.
\end{itemize}

\section*{Animal Handling and Minimization of Pain}
Frogs were located via acoustic searching and visual encounter surveys during diurnal calling windows. Individuals were captured by hand using single-use nitrile gloves to prevent cross-contamination of batrachochytrium dendrobatidis (Bd) pathogens. Animals maintained temporarily for acoustic recording and photomacrography were housed in ventilated terraria supplied with clean moist sphagnum moss and fed flightless fruit flies (\taxa{Drosophila melanogaster}) \textit{ad libitum}.

\section*{Euthanasia and Voucher Preservation}
When terminal voucher preservation was required for osteological and morphological diagnosis, euthanasia was administered via topical application of a 20\% benzocaine gel to the ventral skin, followed by immersion in a buffered 0.2\% tricaine methanesulfonate (MS-222) solution neutralized with sodium bicarbonate (\ce{NaHCO3}, pH 7.4), adhering to the American Veterinary Medical Association (AVMA) Guidelines for the Euthanasia of Animals. Liver and muscle tissue biopsies were harvested immediately post-mortem and preserved in 96\% molecular grade ethanol for genomic analysis. Formalin-fixed vouchers were subsequently transferred to 70\% ethanol for permanent institutional curation at MUSM and QCAZ.
